% Generated draft: requires mathematical and source review.
\documentclass[11pt]{article}
\usepackage[margin=25mm]{geometry}
\usepackage{fontspec}
\setmainfont{Latin Modern Roman}
\usepackage{amsmath,amssymb,mathtools,unicode-math}
\setmathfont{Latin Modern Math}
\usepackage{enumitem,needspace,xurl,hyperref}
\setlength{\parindent}{0pt}
\setlength{\parskip}{5pt}
\setlength{\emergencystretch}{3em}
\newcommand{\N}{\mathbb{N}}
\newcommand{\Z}{\mathbb{Z}}
\newcommand{\Q}{\mathbb{Q}}
\newcommand{\R}{\mathbb{R}}
\newcommand{\C}{\mathbb{C}}
\newcommand{\F}{\mathbb{F}}
\newcommand{\A}{\mathbb{A}}
\newcommand{\PP}{\mathbb{P}}
\newcommand{\E}{\mathbb{E}}
\newcommand{\eps}{\varepsilon}
\newcommand{\dd}{\,\mathrm{d}}
\begin{document}
\section*{UnsolvedMath ID: 700431}
\textbf{Source problem code:} catalogue entry
\subsection{Definitions and mathematical statement}
For a finite group \(H\), define its \(2\)-rank by
\[
m_2(H)=\max\{d\in\Z_{\geq 0}: H\text{ contains a subgroup isomorphic to }(\Z/2\Z)^d\}.
\]
For \(x\in G\), let
\[
C_G(x)=\{g\in G:gx=xg\}
\]
be its centralizer. Interpreting the source's term ``\(3\)-element'' as a nonidentity element of order a power of \(3\), the problem is to describe, up to isomorphism, all finite groups \(G\) such that
\[
\forall x\in G,\qquad
\bigl(|x|=3^k\text{ for some integer }k\geq1\bigr)
\Longrightarrow m_2(C_G(x))\leq2.
\]
Equivalently, it suffices to require this inequality for every element of order exactly \(3\): if \(|x|=3^k\), then \(x^{3^{k-1}}\) has order \(3\) and \(C_G(x)\subseteq C_G(x^{3^{k-1}})\). Groups with no nonidentity \(3\)-elements satisfy the condition vacuously. No simplicity or solvability hypothesis is imposed.

The source does not explicitly exclude the identity from ``\(3\)-elements.'' This draft excludes it; if the identity were included, the condition would instead be equivalent to \(m_2(G)\leq2\), since \(C_G(1)=G\).
\subsection{Short English statement}
Describe all finite groups in which the centralizer of every nonidentity element of \(3\)-power order has \(2\)-rank at most \(2\), equivalently contains no elementary abelian subgroup of order \(8\). The supplied motivation is that many sporadic finite simple groups have this property; the requested classification concerns all finite groups.
\subsection{Sources}
\begin{itemize}
\item Local UnsolvedMath registry, permanent ID \(700431\), exact source problem code \texttt{catalogue entry}, title ``Kourovka Notebook Problem 7.34,'' supplied upstream revision \texttt{d371ee917ac9647fbb06cd6fc529deec6f3dda82}.
\item Supplied source attribution: \emph{The Kourovka Notebook (Unsolved Problems in Group Theory)}, problem 7.34, 7th issue (1980), as recorded in the 21st issue edition (2026), arXiv:1401.0300v48. Proposed by A. A. Makhn\"ev in 1980. The supplied collection description names E. I. Khukhro and V. D. Mazurov as editors.
\item The registry labels the problem ``open.'' This is a supplied status label, not an independently verified assessment. The cited source and its edition have not been independently inspected for this draft.
\end{itemize}
\end{document}
